\documentclass[book.tex]{subfiles}

\begin{document}
\chapter{Calculus}
\label{chap:calculus}
\noindent\rule{11cm}{0.4pt} \\
\textbf{Prerequisites:} \autoref{chap:math_basics} \\
\textbf{Difficulty Level:} * \\
\noindent\rule{11cm}{0.4pt}
\vspace{5mm}

This chapter introduces the basics of calculus, starting from continuity and limits and builds up to Taylor series and antiderivatives.

\section{Continuity and Limits}

In colloquial language, a continuous function is one in which there are no abrupt shift in values. Adding a little bit more formality, suppose that we have a function $y = f(x)$ (read as "f of x"). Let us suppose that $\epsilon$ (read as "epsilon") is an "infinitesimally small" quantity. Then we say that $f$ is a continuous function if the shift $f(x + \epsilon) - f(x)$ is also an infinitesimally small quantity. In the language of "limits", we would write this relationship as 

\begin{equation*}
    \lim_{\epsilon \to 0} f(x + \epsilon) = f(x)
\end{equation*}

(Read this as "limit as epsilon approaches zero of f of x plus epsilon equals f of x".) Many of you may already be familiar with the notion of a limit from past mathematics classes. But for those of you who aren't familiar, let's try to provide a little more intuition. The idea is that the quantity $\epsilon$ needs to "get close" to zero and as it does so, $f(x + \epsilon)$ should "get close" to $f(x)$. The underlying concept here is that of a \textit{metric}, where a \textit{metric} is a function $d: A \times A \to [0, \infty)$ that satisfies the following 3 axioms:

\begin{enumerate}
    \item $d(x, y) = 0$ if and only if $x = y$
    \item $d(x, y) = d(y, x)$ for all $x, y$
    \item $d(x, y) \leq d(x, z) + d(z, y)$ for all $z$
\end{enumerate}

On the real numbers $\mathbb{R}$ (which we haven't yet formally defined), we define the metric as the absolute value of the difference:
\[ d(x, y) = |x - y| \]

A set $A$ with a metric $d$ defined is said to be a "metric space." With this definition in place, we can define the continuity with the "$(\epsilon, \delta)$" definition. That is, a function $f$ on a metric space $A$ is continuous if for every $\epsilon > 0$ there exists a $\delta$ such that if $d(x, y) < \delta$ then $d(f(x), f(y)) < \epsilon$.

This definition looks a little formal, but the idea is simply that as $x$ "gets close" to $y$, then $f(x)$ must "get close" to $f(y)$.


\section{Differential Calculus}

In this section, we'll do a rapid introduction to the calculus that we'll use throughout this book. Let $f: \mathbb{R} \to \mathbb{R}$ be a function. The derivative of $f$, is a function $\frac{df}{dx}: \mathbb{R} \to \mathbb{R}$ which is defined by

\begin{align}
\frac{df}{dx} = \lim_{\Delta x \to 0} \frac{f(x + \Delta x) - f(x)}{\Delta x}
\end{align}

The quantity $\frac{df}{dx}$ is read as "dee f dee x" and is a measure of the rate of change of the function $f$ with respect to the input $x$. We also often use an alternative notation for derivatives

\begin{align}
    f' &= \frac{df}{dx} 
\end{align}
which we read as "f prime". Let's work through a simple example. We will compute $\frac{df}{dx}$ for the following function
\begin{align}
    f(x) = 5 x + 6
\end{align}
We use the definition to write
\begin{align}
    \frac{df}{dx} &= \lim_{\Delta x \to 0} \frac{f(x+\Delta x) - f(x)}{\Delta x} \\
    &=  \lim_{\Delta x \to 0} \frac{5(x+\Delta x) + 6 - (5x + 6)}{\Delta x} \\
    &=  \lim_{\Delta x \to 0} \frac{5\Delta x }{\Delta x} \\
    &= 5
\end{align}
That is, $\frac{df}{dx}$ is the slope of $f$. To compute derivatives for more complex functions, you can use the formal definition, but it is often simpler to use some shortcuts that mathematicians have discovered. The most powerful such shortcut is the chain rule which states
\begin{align}
    \frac{d f(g(x))}{dx} &= \frac{df}{dx}(g(x)) \frac{dg}{dx}
\end{align}
We will use the chain rule systematically later in the book to compute the derivatives of very complex functions. We will also note that we will have reason to chain derivatives. Consider the quantity
\begin{align}
    \frac{d^2 f}{dx^2}
\end{align}
read as "d square f, d x square", which denotes the derivative with respect to $x$ of $\frac{df}{dx}$. We can also define higher derivatives such as $\frac{d^3f}{dx^3}$ and so on. 

Note that this function $\frac{df}{dx}$ doesn't always exist. For example, consider the function $g(x) = |x|$ \\

\begin{marginfigure}
   \includegraphics[width=\linewidth]{figures/Mathematical Foundations/calculus/absolute_value_x.png}
   \caption{A plot of $f(x) = |x|$}
\label{fig:abs_x}
\end{marginfigure}
% TODO: This was generated with wolfram apha. Is that OK to use in the book?

Let's try evaluating the definition at $x = 0$. We get

\begin{align}
    \frac{dg}{dx} = \lim_{\Delta x \to 0} \frac{|x + \Delta x| - |x|}{\Delta x}
\end{align}

There's a bit of subtlety here. Note that if we constrain $\Delta x > 0$, we get $\frac{dg}{dx} = 1$ and if we constrain $\Delta x < 0$, we get $\frac{dg}{dx} = -1$! In practice, what's done numerically is that we handle points such as $0$ by picking one of these derivatives. (Formally, all the numbers in the range $[-1, 1]$ are called \textit{subderivatives} of $|x|$). We say that a function is differentiable if it has a unique value for this limit at all points. A function such as $|x|$ is said to be almost-everywhere differentiable.

\subsection{Taylor Series}

A Taylor series expands a function out in terms of its derivatives. Here is an example of a Taylor expansion around $0$.

\begin{align}
    f(x) = f(0) + f'(0)x + \frac{f''(0)^2}{2!}x^2 + \dotsc
\end{align}

We can define the Taylor series about any point $a$ in fact as
\begin{align}
    f(a) &= f(a) + f'(x-a)(x-a) + \frac{f''(x-a)}{2!} (x-a)^2 + \dotsc
\end{align}

Taylor series are fundamental tools in physics that allow for approximation of a complex function around a local point. Note that Taylor series only converge in a region of the approximation point and not globally typically. One important counterexample to this rule of thumb is the Taylor series for $e^x$ which converges globally

\begin{align}
    e^x = 1 + x + \frac{x^2}{2!} + \frac{x^3}{3!} + \dotsc
\end{align}

\subsection{Multivariate Differentiation}

Suppose that we have a function of two variables
\begin{align}
    f(x, y) &= x^2 + y^2
\end{align}
To find the rate of change of $f$ with respect to either argument, we use partial derivatives defined as
\begin{align}
    \frac{\partial f}{\partial x} &= \lim_{\Delta x \to 0} \frac{f(x + \Delta x, y) - f(x, y)}{\Delta x }\\
    \frac{\partial f}{\partial y} &= \lim_{\Delta y \to 0} \frac{f(x, y + \Delta y) - f(x, y)}{\Delta y}
\end{align}

Read $\frac{\partial f}{\partial x}$ as "partial f partial x" or alternatively as "doh f doh x". The idea is that we hold the other variable constant and take the ordinary derivative with respect to only one argument.

Let $f: \mathbb{R}^N \to \mathbb{R}$ be a differentiable function. We can define partial derivatives $\frac{\partial f}{\partial x_i}$ with respect to the different arguments, but is there a way to define a single joint derivative of $f$? There are many variables here, so which direction do we take the derivative from? The usual solution is to define a directional derivative in a direction $\vec{u}$. We define the directional derivative as

\begin{align}
    \nabla_{\vec{u}}(f) = \lim_{h \to 0} \frac{f(x + h\vec{u}) - f(x)}{h}
\end{align}

It's possible to prove (left as an exercise) that

\begin{align}\nabla_{\vec{u}}(f) = \frac{df}{dx_1}u_1 + \cdots + \frac{df}{dx_N} u_N \end{align}

We define the gradient $\nabla f$ as follows

\begin{align} \nabla f = \begin{bmatrix}
    \frac{df}{dx_1} \\
    \vdots \\
    \frac{df}{dx_N} \\
    \end{bmatrix} \end{align}
    
Note that then the directional derivative is then simply 
\begin{align}\nabla_{\vec{u}}(f) = (\nabla f)^T \vec{u}\end{align}
Note that the gradient is the direction in which the directional derivative is maximized (proof left as an exercise).

Let's now take one more step to generalize. Let $f: \mathbb{R}^N \to \mathbb{R}^M$. What is the derivative of $f$?  It turns out that there are a number of reasonable answers, but we will often find ourselves working with the \textit{Jacobian}. The Jacobian $\mathcal{J}$ is defined as 

\begin{align} \mathcal{J}(f) = \begin{bmatrix}
    \frac{\partial f_1}{dx_1} & \cdots & \frac{\partial f_1}{dx_N} \\
    & \vdots & \\
    \frac{\partial f_M}{x_1} & \cdots & \frac{\partial f_M}{x_N} \\
    \end{bmatrix}
\end{align}
Recall that $\frac{\partial f_i}{\partial x_i}$ is the partial derivative of function $f_i$ with respect to variable $x_i$.
Here we write $f = (f_1, \dotsc, f_M)$ as $M$ functions of type $\mathbb{R}^N \to \mathbb{R}$. Note that $\mathcal{J}(f) \in \mathbb{R}^{M \times N}$.

\section{An introduction to integral calculus}

Integration is the opposite of differentiation. Let $f: \mathbb{R} \to \mathbb{R}$ be a function. We denote the antiderivative of a function by the notation
\begin{equation}
    \int f dx
\end{equation}
The antiderivative is defined by its relation to the derivative. Broadly speaking, the antiderivative $\int f dx$ is a function $\mathbb{R} \to \mathbb{R}$ such that
\begin{equation}
    \frac{d}{dx} \int f dx = f
    \label{antidereqn}
\end{equation}
Antiderivatives are not unique. The reason for this is the fact that $\frac{d}{dx}(C) = 0$ where $C \in \mathbb{R}$ is any constant. This means that the antiderivative is only defined up to a constant. We'll have you work through a few such definitions in the exercises.

It's possible to compute antiderivatives for simple functions directly from the defining Equation \ref{antidereqn}. 

\subsection{Exercises}

\begin{enumerate}
\item Prove that the derivative of $f(x) = e^x$ is $e^x$. Use the fact that
\begin{align}
    e^x &= 1 + x + \frac{x^2}{2!} + \frac{x^3}{3!} + \dotsc
\end{align}
\item Prove that the antiderivative of $f(x) = x^n$ is $\frac{x^{n+1}}{n+1} + C$ by using the defining Equation \ref{antidereqn}.
\item Prove that the antiderivative of $f(x) = e^x$ is $e^x + C$ by using the definiting Equation \ref{antidereqn}. Hint: Consider the Taylor series.
\end{enumerate}



\end{document}